% ============================================================
% RESUMEN PARA EL DIRECTOR DE TESIS
% Eficiencia hospitalaria, estructura organizativa y teoría de agencia
% José María Elena Izquierdo · Universidad de Salamanca · 2026
% ============================================================

\documentclass[12pt, a4paper]{article}

% ── Codificación y lengua ───────────────────────────────────
\usepackage[utf8]{inputenc}
\usepackage[T1]{fontenc}
\usepackage[spanish,es-nodecimaldot]{babel}

% ── Tipografía ──────────────────────────────────────────────
\usepackage{lmodern}
\usepackage{microtype}

% ── Matemáticas ─────────────────────────────────────────────
\usepackage{amsmath, amssymb, bm}

% ── Diseño de página ────────────────────────────────────────
\usepackage{geometry}
\geometry{top=2.5cm, bottom=2.5cm, left=3cm, right=2.5cm}
\usepackage{setspace}
\onehalfspacing

% ── Tablas ──────────────────────────────────────────────────
\usepackage{booktabs}
\usepackage{array}
\usepackage{multirow}
\usepackage{tabularx}
\usepackage{threeparttable}
\usepackage{longtable}

% ── Color y cajas ───────────────────────────────────────────
\usepackage[dvipsnames,svgnames,x11names]{xcolor}
\usepackage{tcolorbox}
\tcbuselibrary{skins, breakable}
\usepackage{mdframed}

% ── Encabezados y pies ──────────────────────────────────────
\usepackage{fancyhdr}
\pagestyle{fancy}
\fancyhf{}
\fancyhead[L]{\small\itshape Eficiencia hospitalaria, estructura organizativa y teoría de agencia}
\fancyhead[R]{\small\itshape Documento de trabajo}
\fancyfoot[C]{\thepage}
\renewcommand{\headrulewidth}{0.4pt}

% ── Secciones ───────────────────────────────────────────────
\usepackage{titlesec}
\definecolor{azultesis}{RGB}{31,78,121}
\definecolor{azulclaro}{RGB}{46,109,164}
\definecolor{gristexto}{RGB}{89,89,89}
\definecolor{fondocaja}{RGB}{235,243,251}
\definecolor{fondotabla}{RGB}{214,228,240}

\titleformat{\section}
  {\color{azultesis}\large\bfseries}
  {\thesection.}{0.5em}{}
  [\vspace{-2pt}\color{azultesis}\rule{\linewidth}{0.8pt}]

\titleformat{\subsection}
  {\color{azulclaro}\normalsize\bfseries}
  {\thesubsection.}{0.5em}{}

\titlespacing{\section}{0pt}{18pt}{8pt}
\titlespacing{\subsection}{0pt}{12pt}{4pt}

% ── Otros ───────────────────────────────────────────────────
\usepackage{enumitem}
\setlist[itemize]{leftmargin=1.4em, itemsep=2pt, parsep=0pt}
\usepackage{parskip}
\setlength{\parskip}{5pt}
\setlength{\parindent}{0pt}

% ── Hyperref ────────────────────────────────────────────────
\usepackage[colorlinks=true,
  linkcolor={azultesis},
  citecolor={azultesis},
  urlcolor={azultesis}]{hyperref}

% ── Caja de hallazgo central ────────────────────────────────
\newtcolorbox{cajadestacada}{
  enhanced,
  colback=fondocaja,
  colframe=azultesis,
  boxrule=0.5pt,
  leftrule=3pt,
  arc=2pt,
  left=8pt, right=8pt, top=6pt, bottom=6pt,
  fontupper=\small\itshape\color{azultesis}
}

% ── Comandos propios ────────────────────────────────────────
\newcommand{\E}{\mathbb{E}}
\newcommand{\bsym}[1]{\boldsymbol{#1}}
\newcommand{\sig}[1]{\ensuremath{^{#1}}}

% ============================================================
\begin{document}
% ============================================================

% ── PORTADA ─────────────────────────────────────────────────
\thispagestyle{empty}
\vspace*{3cm}

\begin{center}
  {\small\color{gristexto}\MakeUppercase{Documento de trabajo para el director de tesis}}\\[6pt]
  \rule{\linewidth}{1.5pt}\\[18pt]

  {\LARGE\bfseries\color{azultesis}
    EFICIENCIA HOSPITALARIA, ESTRUCTURA ORGANIZATIVA\\[4pt]
    Y TEORÍA DE AGENCIA}\\[10pt]

  {\large\itshape\color{gristexto}
    Un análisis de frontera estocástica del sistema hospitalario español}\\[30pt]

  {\large\bfseries José María Elena Izquierdo}\\[6pt]
  {\normalsize\color{gristexto}
    Departamento de Economía Aplicada · Universidad de Salamanca}\\[6pt]
  {\normalsize\color{gristexto} Marzo 2026}\\[40pt]

  {\small\itshape\color{gristexto} Resumen ejecutivo · ca.\ 10 páginas}
\end{center}

\newpage
\setcounter{page}{1}

% ============================================================
\section{Objetivo y motivación}
% ============================================================

Esta tesis analiza empíricamente cómo la estructura organizativa de los
hospitales (centralizada frente a descentralizada) afecta a la eficiencia
técnica en dos dimensiones diferenciadas del output hospitalario: la
\textbf{cantidad} de actividad (volumen de altas ponderadas por case-mix) y la
\textbf{intensidad} diagnóstico-terapéutica por episodio. El punto de partida
es un modelo de agencia multitarea con derivación formal  que
genera predicciones falsificables y asimétricas sobre ambas dimensiones. La
contribución central es contrastar esas predicciones con datos de panel del
sistema hospitalario español (2010--2023) mediante análisis de frontera
estocástica.

\begin{cajadestacada}
\textbf{Pregunta de investigación:} ¿Producen los hospitales descentralizados
(concesiones, fundaciones, privados concertados) mayor eficiencia en cantidad
pero menor eficiencia en intensidad que los hospitales públicos del SNS? ¿Y
depende esa asimetría del tipo de pago y del grado de complementariedad entre
las tareas del médico?
\end{cajadestacada}

La motivación es doble. Teóricamente, la literatura de agencia multitarea
(Holmström y Milgrom, 1991) predice que los incentivos sobre tareas medibles
desplazan esfuerzo desde las tareas menos verificables, pero esta predicción no
ha sido contrastada empíricamente en el sector hospitalario con un diseño de
identificación riguroso. Empíricamente, el sistema hospitalario español ofrece
una variación institucional relevante: conviven hospitales públicos integrados
del SNS con concesiones administrativas (modelo Alzira), fundaciones públicas y
hospitales privados concertados, con formas de pago distintas y distintos
grados de autonomía en el diseño de contratos internos.

% ============================================================
\section{El modelo teórico de agencia multitarea}
% ============================================================

\subsection{Estructura del modelo}

El modelo describe un juego de tres agentes: un financiador público
(principal), el hospital (agente 1) y el médico (agente 2). Los tres agentes
producen conjuntamente dos outputs con tecnologías de producción asimétricas:

\begin{itemize}
  \item \textbf{$q$ (cantidad):} número de altas ponderadas por case-mix. Es el
    resultado del esfuerzo conjunto del hospital ($a$) y del médico ($t_1$),
    más un shock aleatorio: $q = a + t_1 + \varepsilon_1$.
  \item \textbf{$i$ (intensidad):} nivel de recursos diagnóstico-terapéuticos
    aplicados por episodio. Depende exclusivamente del esfuerzo del médico
    ($t_2$): $i = t_2 + \varepsilon_2$.
\end{itemize}

Esta asimetría en la producción es el núcleo estructural del modelo. El
hospital controla la capacidad instalada y el flujo de pacientes; el médico es
el único agente con conocimiento privado sobre la idoneidad y la intensidad del
tratamiento. Si ambos outputs fueran producidos por los mismos agentes con los
mismos incentivos, las predicciones serían simétricas y no habría hipótesis
falsificable sobre el efecto diferencial de la estructura organizativa.

Los shocks aleatorios $\varepsilon_1$ y $\varepsilon_2$ se distribuyen
normalmente con varianzas $\sigma_1^2$ y $\sigma_2^2$ respectivamente. El
parámetro central del modelo es $\psi_{12}$, que mide la complementariedad o
sustitución entre las dos tareas del médico: si $\psi_{12}>0$, las tareas son
\textit{sustitutas} (más tiempo dedicado a ver pacientes reduce la intensidad
por paciente, caso típico de servicios médicos no quirúrgicos); si
$\psi_{12}<0$, son \textit{complementarias} (el acto de intervenir genera
\textit{per se} mayor intensidad, caso típico de servicios quirúrgicos).

\subsection{Estructuras organizativas y contratos óptimos}

El modelo caracteriza dos estructuras organizativas que difieren en la
arquitectura contractual:

\begin{itemize}
  \item \textbf{Centralizada (hospitales SNS):} el financiador contrata
    simultáneamente con el hospital y con el médico. El contrato médico óptimo
    incluye pagos variables sobre cantidad ($B^C$) e intensidad ($D^C$), cuyos
    valores dependen de $\psi_{12}$, $\sigma_1^2$ y $\sigma_2^2$ y de los
    coeficientes de aversión al riesgo.
  \item \textbf{Descentralizada (concesiones, fundaciones, privados
    concertados):} el financiador contrata solo con el hospital para
    determinados niveles de $q$ e $i$, y es el hospital quien diseña
    autónomamente el contrato con el médico. El hospital internaliza el coste
    del diseño del contrato interno, lo que genera incentivos de mayor potencia
    sobre cantidad pero de menor potencia sobre intensidad cuando las tareas son
    sustitutas.
\end{itemize}

Un resultado central del modelo ---denominado \textit{potencia secuencialmente
inducida}--- es que los incentivos del hospital dependen de $\psi_{12}$ en la
estructura descentralizada pero no en la centralizada: cuando las tareas del
médico se aproximan a la sustitución perfecta, el financiador reduce los
incentivos por cantidad al hospital porque anticipa el efecto de desplazamiento
de esfuerzo sobre la intensidad médica.

Las simulaciones numéricas del modelo (Anexo A de la tesis) muestran que el
resultado cualitativo es robusto a la variación de $\psi_{12}$ en todo el
espectro de complementariedad a sustitución. La descentralización mejora
consistentemente el desempeño en cantidad excepto para valores muy negativos de
$\psi_{12}$ (complementariedad extrema), y deteriora el desempeño en intensidad
cuando las tareas son suficientemente sustitutas.

\subsection{Predicciones contrastables}

Del modelo se derivan cinco predicciones empíricas con anclaje directo en los
parámetros de la ecuación de ineficiencia:

\begin{table}[htbp]
\centering
\caption{Predicciones contrastables del modelo teórico}
\label{tab:predicciones}
\begin{threeparttable}
\begin{tabularx}{\linewidth}{cXcc}
\toprule
\textbf{Pred.} & \textbf{Descripción} & \textbf{Parámetro} &
\textbf{Signo} \\
\midrule
P1 & Descentralización $\downarrow$ ineficiencia en cantidad ($q$)
   & $\alpha_1$ en $\mu^q_{it}$ & $\alpha_1 < 0$ \\[3pt]
P2 & Descentralización $\uparrow$ ineficiencia en intensidad ($i$)
   & $\delta_1$ en $\mu^i_{it}$ & $\delta_1 > 0$ \\[3pt]
P3 & El mecanismo de pago tiene efectos asimétricos entre dimensiones
   & $\alpha_2 \neq \delta_2$ & Signos opuestos \\[3pt]
P4 & $\psi_{12}$ modera asimétricamente el diferencial entre estructuras
   & $\alpha_4 \neq \delta_4$ & Signos opuestos \\[3pt]
P$^*$ & Contraste central: el vector completo de determinantes
         difiere entre dimensiones
   & $\bsym{\alpha} \neq \bsym{\delta}$ & Wald rechazado \\
\bottomrule
\end{tabularx}
\begin{tablenotes}
  \small
  \item $\mu^q_{it}$ y $\mu^i_{it}$ son las medias heterogéneas de los
    términos de ineficiencia en las fronteras de cantidad e intensidad.
    $\bsym{\alpha}$ y $\bsym{\delta}$ son los vectores completos de
    coeficientes de la ecuación de ineficiencia en cada frontera.
\end{tablenotes}
\end{threeparttable}
\end{table}

La predicción P$^*$ es la que mayor exigencia estadística impone: implica que
el vector completo de determinantes institucionales opera de forma
cualitativamente distinta sobre las dos dimensiones del output. Un único término
de ineficiencia $u_{it}$ cancela algebraicamente los efectos opuestos de P1 y
P2, haciendo imposible contrastar P$^*$ con los modelos SFA estándar de un solo
output. Esta observación motiva el diseño empírico en sistema de dos fronteras.

% ============================================================
\section{Estrategia empírica: análisis de frontera estocástica}
% ============================================================

\subsection{Marco general SFA}

El análisis de frontera estocástica descompone la desviación de cada hospital
respecto a la frontera de producción en dos componentes: ruido aleatorio
simétrico ($v_{it}$) e ineficiencia técnica no negativa ($u_{it}$). La
eficiencia técnica individual se recupera mediante el estimador de esperanza
condicional de Jondrow et al.\ (1982):
\[
  TE_{it} = \E\!\left[e^{-u_{it}} \mid \varepsilon_{it}\right]
\]
Para incorporar los determinantes institucionales de la ineficiencia, se adopta
la especificación de Battese y Coelli (1995), donde $u_{it}$ sigue una
distribución normal truncada en cero con media heterogénea $\mu_{it} =
\mathbf{z}_{it}'\bsym{\delta}$.

La frontera de producción se especifica como una función \textbf{translog}
centrada en las medias muestrales, con inputs trabajo (personal médico EJC),
capital físico (camas en funcionamiento), capital tecnológico (índice de
equipamiento diagnóstico-terapéutico), tendencia temporal cuadrática (progreso
técnico) y dummies de cluster de complejidad hospitalaria (cinco grupos). La
especificación translog es estadísticamente preferida sobre Cobb-Douglas en
todos los modelos, con estadísticos LR entre 247 y 594 ($p<0.001$).

\subsection{Ecuación de ineficiencia}

La media heterogénea del término de ineficiencia es la misma en los cuatro
diseños:
\begin{align*}
  \mu_{kit} = \delta_0
    &+ \delta_1\, D^{desc}_{it}
     + \delta_2\, pct^{SNS}_{it}
     + \delta_3\, (D^{desc}_{it} \times pct^{SNS}_{it}) \\
    &+ \delta_4\, ShareQ_{it}
     + \delta_5\, (D^{desc}_{it} \times ShareQ_{it})
     + \sum_{c}\lambda_c\, D^{ccaa}_{c,it}
\end{align*}

Cada variable tiene un anclaje directo en el modelo teórico:
\begin{itemize}
  \item $D^{desc}_{it}$: estructura organizativa (1 = descentralizado, 0 = SNS
    centralizado). Sus coeficientes en las dos fronteras son $\alpha_1$ y
    $\delta_1$ de P1 y P2.
  \item $pct^{SNS}_{it}$: proporción de altas financiadas por el SNS como
    proxy del mecanismo de pago. Capta P3.
  \item $D^{desc} \times pct^{SNS}$: interacción estructura--pago (P3).
  \item $ShareQ_{it}$: proporción de altas quirúrgicas, proxy continua de
    $\psi_{12}$. Capta P4.
  \item $D^{desc} \times ShareQ$: moderación de $\psi_{12}$ sobre el efecto
    diferencial entre estructuras (P4).
\end{itemize}

\subsection{Los cuatro diseños de estimación}

\begin{table}[htbp]
\centering
\caption{Diseños de estimación y su papel en el análisis}
\label{tab:disenos}
\begin{threeparttable}
\small
\begin{tabular}{cp{4.5cm}p{2.8cm}p{4cm}}
\toprule
\textbf{Dis.} & \textbf{Descripción} & \textbf{Hipótesis} &
\textbf{Papel} \\
\midrule
D & Función de distancia output (ODF), un único $u_{it}$
  & P3, P4
  & Punto de partida metodológico \\[4pt]
A & Sistema de dos fronteras: cantidad total vs.\ intensidad diagnóstica ($i_{diag}$)
  & P1, P2, P3, P4, P$^*$
  & Contraste directo del modelo teórico \\[4pt]
B & Sistema de dos fronteras: altas quirúrgicas vs.\ altas médicas
  & P1, P3, P4, P$^*$
  & Contraste con proxy estructural de $\psi_{12}$ \\[4pt]
C & Panel hospital $\times$ servicio $\times$ año, interacción triple
  & P4 (variación intra-hospital)
  & Robustez dentro del hospital \\
\bottomrule
\end{tabular}
\begin{tablenotes}
  \small
  \item Los Diseños A y B son los centrales para el contraste del modelo. El
    Diseño B usa la clasificación quirúrgico/médico como proxy estructural de
    $\psi_{12}$, exógeno a nivel de servicio, evitando la endogeneidad de un
    proxy construido a nivel hospital-año.
\end{tablenotes}
\end{threeparttable}
\end{table}

Los contrastes formales son: (i)~test $t$ unilateral sobre $\alpha_1$ y
$\delta_1$ (P1 y P2 individualmente); (ii)~test $z$ sobre la diferencia
$\alpha_1 - \delta_1$ (contraste de asimetría); (iii)~test de Wald bootstrap
con $B=4{,}999$ repeticiones sobre $H_0: \bsym{\alpha} = \bsym{\delta}$
(predicción P$^*$); (iv)~test LR Kodde-Palm para la existencia de ineficiencia;
y (v)~test LR para la forma funcional (Cobb-Douglas vs.\ translog). La muestra
principal del Diseño A consta de 4.765 observaciones hospital-año (880
hospitales únicos, período 2010--2023).

% ============================================================
\section{Resultados principales}
% ============================================================

\subsection{Síntesis}

\begin{table}[htbp]
\centering
\caption{Síntesis de resultados por predicción}
\label{tab:resultados}
\begin{threeparttable}
\small
\begin{tabular}{cp{3.2cm}p{3cm}p{2.8cm}c}
\toprule
\textbf{Pred.} & \textbf{Diseño A} & \textbf{Veredicto A}
  & \textbf{Diseño B} & \textbf{Veredicto final} \\
\midrule
P1 & $\hat{\alpha}_1=-2.530^{***}$ ($t=-6.98$)
   & Confirmada (7/7 esp.)
   & Confirmada vía $Priv\_Merc$
   & \textbf{Confirmada} \\[4pt]
P2 & $z_{dif}=-2.09^{*}$ (asimetría sig., sin inversión)
   & Parcial
   & $\hat{\delta}_{Merc}<0$
   & \textbf{Parcial} \\[4pt]
P3 & $z_{dif}=2.35^{*}$
   & Confirmada
   & $z_{dif}=5.32^{***}$
   & \textbf{Confirmada} \\[4pt]
P4 & $z_{dif}=-0.84$ (n.s.)
   & No confirmada en A
   & $z_{dif}=-9.89^{***}$
   & \textbf{Confirmada (B)} \\[4pt]
P$^*$ & $LR=10{,}491^{***}$
   & Confirmada
   & $LR=2{,}508^{***}$
   & \textbf{Confirmada} \\
\bottomrule
\end{tabular}
\begin{tablenotes}
  \small
  \item $^{*}p<0.05$; $^{***}p<0.001$. La eficiencia técnica media oscila
    entre 0.683 (frontera de intensidad, Diseño A) y 0.791 (frontera
    quirúrgica, Diseño B).
\end{tablenotes}
\end{threeparttable}
\end{table}

\subsection{Predicción P1: descentralización reduce ineficiencia en cantidad}

El coeficiente de $D^{desc}$ en la frontera de cantidad es $-2.530$ ($t=-6.98$,
$p<0.001$) en la especificación principal. Los hospitales con mayor autonomía en
el diseño de contratos internos producen más actividad por unidad de input. Este
resultado es extraordinariamente robusto: el signo y la significación se
mantienen en las siete especificaciones alternativas del análisis de robustez,
con $t$-estadísticos entre $-3.3$ y $-11.0$.

\subsection{Predicción P2: descentralización aumenta ineficiencia en intensidad}

El coeficiente de $D^{desc}$ en la frontera de intensidad es $-1.477$
($t=-4.21$), negativo en todas las especificaciones, no positivo como predice el
modelo bajo el supuesto de sustitución fuerte. Sin embargo, la diferencia entre
los dos coeficientes es estadísticamente significativa ($z_{dif}=-2.09$,
$p=0.037$): la descentralización reduce la ineficiencia en cantidad
significativamente \textit{más} que en intensidad, en la dirección asimétrica
predicha por el modelo.

\begin{cajadestacada}
\textbf{Interpretación de P2:} el signo negativo de $\hat{\delta}_1$ es
consistente con un valor medio de $\psi_{12}$ insuficiente para revertir el
efecto, o bien con la presencia de mecanismos institucionales de control de
calidad (acreditación autonómica, contratos-programa con estándares de
intensidad) que limitan la sustitución de esfuerzo entre tareas. La asimetría
estadísticamente significativa sí confirma la dirección del efecto predicho por
el modelo.
\end{cajadestacada}

\subsection{Predicción P3: el tipo de pago modera asimétricamente}

El coeficiente de $pct^{SNS}$ es negativo y altamente significativo en ambas
fronteras ($-2.473$ en cantidad, $-2.439$ en intensidad): una mayor financiación
pública reduce la ineficiencia en ambas dimensiones. La interacción
$D^{desc}\times pct^{SNS}$ difiere significativamente entre fronteras
($z_{dif}=2.35$ en A, $z_{dif}=5.32$ en B, ambos $p<0.05$): el tipo de pago
modera el efecto de la descentralización de forma cualitativamente distinta
sobre cantidad e intensidad.

\subsection{Predicción P4: $\psi_{12}$ modera asimétricamente el diferencial}

El resultado individual más robusto de todo el análisis. En el Diseño B,
$ShareQ$ tiene signos completamente opuestos en las dos fronteras: $-7.015$
($t=-8.42$) en quirúrgico y $+4.647$ ($t=5.58$) en médico, con
$z_{dif}=-9.89$ ($p<0.001$). Los hospitales con mayor proporción de actividad
quirúrgica operan en un régimen de complementariedad entre tareas ($\psi_{12}<0$):
el acto quirúrgico genera \textit{per se} mayor intensidad diagnóstica pre y
postoperatoria, lo que mejora la eficiencia relativa en la frontera quirúrgica
pero la deteriora en la frontera médica. La correlación negativa entre
ineficiencias en el Diseño B ($r=-0.140$, $p<0.001$) proporciona evidencia
adicional directa del mecanismo de sustitución de esfuerzo entre tareas.

\subsection{Contraste central P$^*$: los vectores de determinantes difieren}

El test de razón de verosimilitud sobre $H_0: \bsym{\alpha} = \bsym{\delta}$
rechaza abrumadoramente la hipótesis nula en ambos diseños ($LR_A=10{,}491$,
$LR_B=2{,}508$, ambos $p<0.001$). Los determinantes institucionales de la
ineficiencia operan de forma cualitativamente distinta sobre las dos dimensiones
del output. Este resultado no es recuperable con ningún modelo SFA estándar de
un único término de ineficiencia.

% ============================================================
\section{Tests de robustez y especificación}
% ============================================================

\subsection{Tests formales de especificación}

\begin{itemize}
  \item \textbf{Test Kodde-Palm ($H_0: \sigma^2_u=0$):} estadísticos LR entre
    474 y 2.240 ($p<0.001$ en todos los modelos). La ineficiencia técnica es
    estadísticamente significativa.
  \item \textbf{Test LR forma funcional (Cobb-Douglas vs.\ translog):}
    estadísticos entre 247 y 594 ($p<0.001$). La translog es preferida en todos
    los diseños.
  \item \textbf{Correlación entre ineficiencias:} baja en valor absoluto
    ($|r|<0.15$ en ambos diseños), lo que justifica la estimación independiente
    de las dos fronteras.
\end{itemize}

\subsection{Análisis de robustez}

\begin{table}[htbp]
\centering
\caption{Robustez: coeficiente de $D^{desc}$ en fronteras de cantidad e
intensidad}
\label{tab:robustez}
\begin{threeparttable}
\small
\begin{tabular}{clrrrrc}
\toprule
& & \multicolumn{2}{c}{$\hat{\alpha}$ (cantidad)}
  & \multicolumn{2}{c}{$\hat{\delta}$ (intensidad)}
  & \\
\cmidrule(lr){3-4}\cmidrule(lr){5-6}
\textbf{Esp.} & \textbf{Modificación} & Coef. & $t$
  & Coef. & $t$ & \textbf{Asim.} \\
\midrule
\textbf{Princ.} & $i_{diag}$, $D^{desc}$ coalesce
  & $-2.530$ & $-6.98$ & $-1.477$ & $-4.21$ & Sí \\
R1 & Intensidad alternativa ($i_{simple}$)
  & $-1.264$ & $-4.62$ & $-0.704$ & $-3.33$ & Sí \\
R2a & $D^{desc}$ fuente SIAE
  & $-2.954$ & $-9.08$ & $-0.414$ & $-1.82$ & Sí \\
R2b & $D^{desc}$ fuente CNH
  & $-6.498$ & $-8.60$ & $-0.865$ & $-1.69$ & Sí \\
R3 & Distribución half-normal
  & $-5.245$ & $-10.19$ & $-2.788$ & $-5.18$ & Sí \\
R4 & Muestra 2010--2019 (sin COVID)
  & $-2.649$ & $-6.05$ & $-2.880$ & $-4.66$ & Sí \\
R5 & Sin hospitales ONG
  & $-6.009$ & $-10.97$ & $-2.730$ & $-5.24$ & Sí \\
\bottomrule
\end{tabular}
\begin{tablenotes}
  \small
  \item $\hat{\alpha}$ negativo y significativo en 7/7 especificaciones.
    $\hat{\delta}$ significativo al 5\% en 5/7 especificaciones. La columna
    Asim.\ indica si la diferencia $\hat{\alpha}-\hat{\delta}$ es
    estadísticamente significativa.
\end{tablenotes}
\end{threeparttable}
\end{table}

Destaca la especificación R4 (muestra 2010--2019): $\hat{\delta}_1 = -2.880$,
más próximo al valor predicho por el modelo que la especificación principal
($-1.477$), lo que sugiere que el período post-COVID introdujo efectos de
recuperación diferencial que atenúan la asimetría observada. La especificación
R2b (fuente CNH, más granular) produce los coeficientes más elevados en valor
absoluto, coherente con que el CNH identifica mejor a los hospitales
genuinamente descentralizados.

Están pendientes de ejecución: (i)~bootstrap Wald con $B=4{,}999$ para el test
P$^*$ formal; (ii)~sistema bivariado con cópula gaussiana (Kumbhakar, 1997) si
el test de independencia de ineficiencias es rechazado; y (iii)~reestimación
con pesos GRD diferenciados por tipo de GRD quirúrgico/médico, condicionada a
la recepción de los microdatos RAE-CMBD del Ministerio de Sanidad.

% ============================================================
\section{Implicaciones de política hospitalaria \texorpdfstring{\normalfont\small\itshape\color{gristexto}[capítulo en elaboración]}{}}
% ============================================================

\begin{mdframed}[
  linecolor=orange!70!black,
  linewidth=0.8pt,
  backgroundcolor=orange!6,
  innerleftmargin=10pt,
  innerrightmargin=10pt,
  innertopmargin=7pt,
  innerbottommargin=7pt,
  skipabove=6pt,
  skipbelow=6pt
]
\small\itshape\color{orange!80!black}
\textbf{Nota:} este capítulo está actualmente en elaboración. Las propuestas
recogidas a continuación representan las líneas de política que se derivan de
forma más directa de los resultados empíricos disponibles, pero el análisis no
está cerrado. En particular, están pendientes de incorporar: la discusión sobre
los límites de la evidencia para fundamentar recomendaciones causales, el
tratamiento diferenciado por tipo de comunidad autónoma y modelo de
financiación, y las implicaciones específicas que se deriven de los resultados
del bootstrap Wald y del sistema bivariado con cópula cuando estén disponibles.
\end{mdframed}

\vspace{4pt}

Los resultados tienen implicaciones directas para el diseño institucional del
sistema hospitalario. La evidencia confirma que la estructura organizativa tiene
efectos estadísticamente significativos y cualitativamente distintos sobre las
dos dimensiones del output. Un enfoque de política de ``talla única'' ---ya sea
centralización total o descentralización generalizada--- ignora esta
multidimensionalidad.

\begin{enumerate}[label=\textbf{\arabic*.}, leftmargin=*]

  \item \textbf{Evaluación multidimensional del rendimiento hospitalario.} Los
    sistemas de evaluación deberían incorporar medidas separadas de eficiencia
    en volumen y en intensidad diagnóstico-terapéutica. Una métrica global agrega
    efectos de signo contrario y puede confundir a los reguladores sobre el
    desempeño real de los hospitales descentralizados. El diseño de esta tesis
    ---sistema de dos fronteras estocásticas--- puede implementarse con los datos
    disponibles del SIAE y del CMBD como herramienta de regulación por
    comparación (\textit{yardstick competition}).

  \item \textbf{Extensión del pago por actividad diferenciado.} El coeficiente
    de $pct^{SNS}$ sugiere que una mayor proporción de financiación por
    actividad reduce la ineficiencia en ambas dimensiones. La extensión del
    sistema de pago por GRD a toda la red concertada, con tarifas que
    diferencien entre servicios quirúrgicos y médicos, es la medida con mayor
    apoyo empírico.

  \item \textbf{Inclusión de indicadores de intensidad en los
    contratos-programa.} Los contratos-programa autonómicos deberían incluir
    indicadores verificables de intensidad diagnóstico-terapéutica (número de
    procedimientos por alta, tasa de diagnóstico por imagen, estancia media
    ajustada por severidad) vinculados a una parte de la financiación, para
    contrarrestar el efecto de sustitución de esfuerzo predicho por el modelo
    cuando $\psi_{12}>0$.

  \item \textbf{Cláusulas de calidad en contratos de concesión.} Para
    hospitales en régimen de concesión administrativa, los pliegos de
    condiciones deberían incluir umbrales mínimos verificables de intensidad
    diagnóstica y auditoría periódica. Los resultados sugieren que los
    mecanismos de control de calidad ya existentes (acreditación autonómica)
    limitan el deterioro de la intensidad pero no lo eliminan.

  \item \textbf{Diseño de reformas de descentralización.} La evidencia es
    consistente con los efectos causales del modelo teórico, pero no los
    establece de forma concluyente (la asignación de hospitales a formas
    descentralizadas no es aleatoria). Para establecer causalidad sería
    necesario explotar variación exógena derivada de cambios normativos
    autonómicos (reversiones de concesiones en la Comunidad Valenciana y
    Madrid).

\end{enumerate}

\begin{cajadestacada}
\textbf{Síntesis de política:} los resultados apuntan a que la clave no es
elegir entre centralización y descentralización, sino diseñar contratos
hospitalarios que aprovechen las ventajas de la descentralización en volumen
(P1 confirmada) mientras se incorporan mecanismos institucionales que
contrarresten su desventaja en intensidad (P2 confirmada en sentido débil). El
tipo de pago y el mix de servicios quirúrgico/médico del hospital son los
moderadores institucionales que el diseño de contratos debería tener en cuenta.
\end{cajadestacada}

% ============================================================
\section{Estado actual y estructura de la tesis}
% ============================================================

La tesis está estructurada en ocho capítulos y tres anexos. El modelo teórico
(Cap.~2) y la estrategia empírica (Cap.~3) están completos. Los capítulos de
datos (4), resultados (5) y robustez (6) recogen los resultados actualmente
disponibles. El capítulo de conclusiones (8) está redactado en borrador. El
capítulo de política hospitalaria (7) está en elaboración: las propuestas
recogidas en la sección anterior son una primera versión que se ampliará y
matizará conforme se completen los análisis de robustez pendientes. La Tabla~\ref{tab:pendientes}
recoge el conjunto de tareas en curso, tanto econométricas como de redacción:

\begin{table}[htbp]
\centering
\caption{Tareas pendientes por orden de prioridad}
\label{tab:pendientes}
\begin{tabular}{p{5.5cm}p{2cm}p{5cm}}
\toprule
\textbf{Tarea} & \textbf{Prioridad} & \textbf{Estado} \\
\midrule
Bootstrap Wald P$^*$ ($B=4{,}999$)
  & Alta & Pendiente ($\sim$6h cómputo) \\[3pt]
Sistema bivariado con cópula gaussiana
  & Alta & Pendiente (depende de test $\rho$) \\[3pt]
Diseño A con pesos GRD diferenciados Q/M
  & Alta & Pendiente (espera RAE-CMBD) \\[3pt]
Cap.~7 --- Política hospitalaria (redacción completa)
  & Alta & En elaboración --- primera versión disponible \\[3pt]
Control dummies COVID en muestra principal
  & Media & Estimación R4 disponible \\[3pt]
Diseño D (ODF): resolver SE=NaN
  & Baja & Orientativo --- baja prioridad \\
\bottomrule
\end{tabular}
\end{table}

\vspace{1.5cm}
\begin{center}
  \rule{0.5\linewidth}{0.4pt}\\[6pt]
  {\small\itshape\color{gristexto}
    Documento de trabajo --- Marzo 2026 --- José María Elena Izquierdo}
\end{center}

\end{document}
